\chapter{What a Differential Equation Is}

In Chapters~32 through~35, we explored the rich geometry of several variables: partial derivatives, gradients, multiple integrals, and the integral theorems of Green, Stokes, and Gauss. Those theorems dealt with static vector fields defined across curves, surfaces, and volumes in space.

Now we pivot from the static geometry of space to the dynamic geometry of time. Why return from multivariable vector fields to functions of a single variable? Because before we can understand how an entire spatial field evolves in time---the domain of partial differential equations like the heat, wave, and Navier-Stokes equations---we must first master how a single quantity moves forward under a dynamic law. And as you will discover in this chapter, vector fields have not been left behind: a differential equation in the plane assigns a direction to every coordinate, and its solution curves are streamlines flowing through that field.

Every chapter until now has handed you a function and asked questions about it. Here is a curve; what is its slope? Here is a region; what is its flux or volume? The function came first, fully formed, and calculus was the lens through which we studied it. In this chapter, we consider the reverse perspective, which reflects the original motivation for inventing calculus.

For in the real world the function almost never comes first. What comes first is a \emph{law} -- a statement about how a quantity is changing right now, expressed in terms of how much of that quantity there is right now. A population grows faster when there are more individuals to reproduce. A hot cup of coffee cools faster when it is further above room temperature. A falling stone speeds up at a rate fixed by gravity. In each case nature does not tell you the function $y(t)$ directly. Nature tells you something about $y$ and its rate of change $\dv{y}{t}$, tangled together in a single equation, and leaves you the job of untangling it -- of finding the actual function that obeys the law.

An equation that ties a function to its own derivatives is called a \emph{differential equation} (in plain terms: a rule relating what a quantity is doing right now to how fast it is changing). Solving one means finding the functions that satisfy the law. This is the central problem of mathematical physics, of population biology, of engineering, of economics -- of nearly every quantitative science. It is, in a real sense, the question calculus was built to answer. Newton developed the derivative in large part to formulate and solve the differential equations of motion governing celestial bodies.

Here we get clear on \emph{what these equations are}, what it means to solve one, why a single law can have a whole family of solutions until an initial condition pins down just one, and how you can \emph{see} the solutions as curves threading through a field of little slope marks, even before you have solved anything algebraically.

\section{A law about change}
\label{ch36:sec-law}

Start with the simplest possible example, built from a sentence rather than from symbols.

\begin{intuitionbox}
\emph{A quantity grows at a rate equal to its current size.} That sentence states a differential equation in words. Piece by piece: ``A quantity'' is some function $y(t)$ -- maybe a population, maybe a bank balance. ``Grows at a rate'' is its derivative $\dv{y}{t}$. ``Equal to its current size'' means that rate equals $y$ itself. Put the three pieces together and the sentence becomes
\[
  \dv{y}{t} = y.
\]
That is it. That is a differential equation. It does not tell you what $y$ is. It tells you a \emph{relationship} that the true $y$ -- whatever it turns out to be -- must satisfy at every instant.
\end{intuitionbox}

Notice what kind of object an equation like $\dv{y}{t} = y$ is. An ordinary algebraic equation, like $x^2 = y$, relates \emph{numbers}: plug in a number for $x$, get a number for $y$. A differential equation relates a \emph{function} to its own derivative. Its unknown is not a number but an entire function. When we ask for a solution, we are asking: \emph{which function, when you differentiate it, gives back a result that fits this equation?}

For $\dv{y}{t} = y$, you may already feel the answer pressing on you. We need a function that is its own derivative. The exponential function $y = e^t$ has this property: $\ddx{e^t}$ -- read with $t$ in place of $x$ -- equals $e^t$ again. So $y = e^t$ solves the equation. And it is not the only one: $y = 3e^t$ works too, because its derivative is $3e^t$, which equals the function itself. So does $y = -\tfrac{1}{2}e^t$, and $y = Ce^t$ for \emph{any} constant $C$. This infinite collection $y = Ce^t$ is called the \emph{general solution} (the full family of formulas that satisfy the differential equation, indexed by an arbitrary constant). When a specific constant is chosen, such as $y = 3e^t$, we call it a \emph{particular solution} (a single concrete trajectory). We will return to this whole family shortly; for now just register the shape of things. One law, many solutions.

\begin{definition}[Differential equation; ordinary and partial]
\label{ch36:def-de}
A \emph{differential equation} is an equation relating an unknown function to one or more of its derivatives. If the unknown function depends on a single independent variable -- so that only ordinary derivatives appear -- the equation is called an \emph{ordinary differential equation}, abbreviated \textbf{ODE}. (An ordinary derivative tracks how a quantity changes along a single parameter, like time.) If the unknown depends on several variables and partial derivatives appear, it is a \emph{partial differential equation}, abbreviated \textbf{PDE}. (A partial derivative measures the slope of a surface in one chosen direction while all other coordinates remain constant.) This chapter, and the three that follow, are about ODEs.
\end{definition}

The distinction in the definition matters. A PDE like the heat equation $\pdv{u}{t} = \frac{\partial^2 u}{\partial x^2}$ involves a function $u(x,t)$ of two variables -- here the temperature $u$ at position $x$ and time $t$ -- and is genuinely harder; PDEs are a subject of their own. By restricting to a single independent variable -- usually time $t$ or position $x$ -- we keep the unknown a curve in the plane, something we can draw. That restriction is exactly what makes the geometric picture in this chapter possible.

\begin{historybox}
The phrase ``differential equation'' goes back to Gottfried Wilhelm Leibniz, who in a 1676 manuscript was already manipulating equations relating a variable quantity to its differential $\dd y$ (his notation for an infinitesimal change). But the subject truly ignited with Isaac Newton. In the \emph{Principia} (1687) and his earlier work on fluxions (Newton's term for rates of change, imagining quantities flowing smoothly through time), Newton's deepest insight was that the laws of nature are most naturally stated as equations of \emph{rates}. His second law of motion, $F = m\,\frac{\dd^2 x}{\dd t^2}$ where $\frac{\dd^2 x}{\dd t^2}$ is the second derivative (acceleration), is a differential equation for position. Newton solved this equation for a planet orbiting the Sun under an inverse-square gravitational pull and mathematically \emph{derived} Kepler's elliptical orbits -- proving that a local law of change reproduces the orbital motions of celestial bodies.

In the decades that followed, the Swiss brothers Jakob and Johann Bernoulli turned differential equations into an active field of investigation. They wrestled with physical curve puzzles, notably the \emph{catenary} (the U-shaped curve naturally assumed by a heavy hanging chain supported at both ends) and the celebrated 1696 \emph{brachistochrone} problem (from the Greek for ``shortest time,'' asking what curve allows a frictionless bead to slide between two points in the least possible time). Armed with early infinitesimal methods, power series, and ingenious coordinate substitutions, the Bernoullis and Leonhard Euler solved these problems (the brachistochrone turned out to be an inverted cycloid).

Yet the era was plagued by controversy. A bitter priority feud raged between Newton and Leibniz over who first discovered calculus. A second, deeper debate arose over what actually counted as a ``solution.'' When an equation could not be solved in terms of elementary functions (familiar algebraic combinations of polynomials, roots, sines, and logarithms), was it truly solved? Euler recognized that differential equations themselves define brand-new transcendental functions that expand mathematics beyond elementary algebra.

These methods formed the mathematical basis of classical physics: celestial navigation, wave propagation, acoustics, and the engineering of the Industrial Revolution (steam engines, bridges, and fluid pumps). What remains open today? Most nonlinear differential equations cannot be solved by explicit formulas. Even the simple-looking gravitational motion of three bodies (the classical \emph{three-body problem}) was shown by Henri Poincar\'e in the 1880s to be non-integrable and prone to deterministic chaos -- extreme sensitivity where tiny shifts in starting positions cause wildly diverging futures. And in partial differential equations, proving whether smooth solutions to the three-dimensional Navier-Stokes equations (which govern fluid flow like air turbulence and ocean currents) always exist for all time or can develop infinite singularities remains a famous million-dollar open problem.
\end{historybox}

\subsection{The general shape}

Most of the equations we meet can be written so that the highest derivative sits alone on the left:
\[
  \dv{y}{t} = f(t, y),
\]
read aloud as ``the rate of change of $y$ is some specified function of the current time $t$ and the current value $y$.'' The right-hand side $f(t,y)$ is the \emph{law}: hand it where you are ($t$) and what you have ($y$), and it returns how fast $y$ is changing. Our first example $\dv{y}{t} = y$ is the case $f(t,y) = y$ -- the rate depends on $y$ but not directly on $t$. The slightly different equation $\dv{y}{t} = t$ has $f(t,y) = t$ -- the rate depends on the clock but not on the current value. Both are legitimate, and they behave very differently, as we will see.

\begin{happeningbox}
A differential equation $\dv{y}{t} = f(t,y)$ is a \emph{rule for the slope of a curve, given the curve's location}. It requires that at any point $(t,y)$ in the plane, a solution curve must have slope $f(t,y)$. A solution is a curve that satisfies this slope condition everywhere. Solving the equation means finding a function whose graph matches these local slopes at every point. This geometric interpretation is formalized by the slope field in Section~\ref{ch36:sec-slopefield}.
\end{happeningbox}

\section{What it means to solve one}
\label{ch36:sec-solve}

Solving a differential equation is unlike solving the equations you are used to. When you solve $x^2 - 5x + 6 = 0$ you are hunting for \emph{numbers} -- here $x=2$ and $x=3$. When you solve a differential equation you are hunting for \emph{functions}. The answer is not a value but a formula, a whole curve.

\begin{definition}[Solution of an ODE]
\label{ch36:def-solution}
A \emph{solution} of the differential equation $\dv{y}{t} = f(t,y)$ on an interval $I$ is a function $y = \varphi(t)$, differentiable on $I$, such that substituting $\varphi$ for $y$ makes the equation true for every $t$ in $I$. That is, $\varphi'(t) = f\big(t, \varphi(t)\big)$ holds identically on $I$. Verifying a proposed solution is therefore a purely mechanical act: differentiate it, substitute, and check that the two sides agree.
\end{definition}

The last sentence of the definition is important: \emph{checking} a proposed solution
is always straightforward, even when \emph{finding} it was difficult. You differentiate, you substitute, you see whether the equation holds. Here is that check on our running example.

\begin{example}[Verifying a solution]
\label{ch36:ex-verify}
Claim: $y = Ce^t$ solves $\dv{y}{t} = y$ for every constant $C$. To check, differentiate the candidate:
\[
  \dv{y}{t} = \frac{\dd}{\dd t}(Ce^t) = Ce^t.
\]
And the right-hand side of the equation is just $y = Ce^t$. The two agree -- $Ce^t = Ce^t$ -- for every $t$ and every $C$. So the claim is verified. We have not \emph{derived} this solution; we have only confirmed it. But confirmation is genuine knowledge: we now know with certainty that every function of the form $Ce^t$ obeys the law.
\end{example}

\begin{example}[A solution by direct integration]
\label{ch36:ex-integrate}
Consider another simple equation, $\dv{y}{t} = t$. Here the right-hand side depends only on $t$. An equation of this form, $\dv{y}{t} = g(t)$, is solved by integration, because it asks for a function whose derivative is $g(t)$---the standard antiderivative problem. So
\[
  y = \int t \,\dd t = \tfrac{1}{2}t^2 + C.
\]
Again a constant of integration appears, and again it is a \emph{family} of solutions, one for each value of $C$. Every parabola $\tfrac{1}{2}t^2 + C$ is a vertical shift of every other, and each one has slope exactly $t$ at every point, as the equation demands.
\end{example}

These two examples reveal something structural that you must internalize: \textbf{a first-order differential equation does not have one solution; it has a one-parameter family of solutions.} The arbitrary constant $C$ is inherent to the problem rather than an artifact of our solution method.

\begin{whybox}
A free constant appears because the differential equation specifies only the \emph{slope} of the solution curve, not its \emph{height}. For $\dv{y}{t} = t$, the curve must have slope $t$ at position $t$, but this condition does not determine the vertical placement. Any solution can be shifted vertically by a constant without changing its slope. This vertical freedom corresponds to the integration constant $C$. In general, a first-order equation determines the solution up to an arbitrary constant, reflecting the initial value.
\end{whybox}

\begin{figure}[H]
\centering
\begin{tikzpicture}
\begin{axis}[
  width=11cm, height=6.8cm,
  xlabel={$t$}, ylabel={$y$},
  xlabel style={anchor=south west},
  xmin=-2.6, xmax=2.6, ymin=-2.2, ymax=3.8,
  axis lines=middle,
  xtick={-2,-1,1,2}, ytick={-2,-1,1,2,3},
  enlargelimits=false,
  clip=true,
]
\addplot[gray!50, thick, samples=60, domain=-2.6:2.6]{0.5*x^2 - 1.5};
\addplot[gray!50, thick, samples=60, domain=-2.6:2.6]{0.5*x^2 - 0.5};
\addplot[curveblue, very thick, samples=80, domain=-2.6:2.6]{0.5*x^2 + 0.5};
\addplot[gray!50, thick, samples=60, domain=-2.6:2.6]{0.5*x^2 + 1.5};

\addplot[structurepurple, dashed, thick, domain=-2.2:3.8] coordinates {(1.2, -2.2) (1.2, 3.8)};
\node[structurepurple, font=\scriptsize\bfseries, anchor=south] at (axis cs:1.2, 3.1) {$t = 1.2$};

\addplot[tangentred, very thick, domain=0.85:1.55]{1.2*(x - 1.2) + (0.5*1.2^2 - 1.5)};
\addplot[tangentred, very thick, domain=0.85:1.55]{1.2*(x - 1.2) + (0.5*1.2^2 - 0.5)};
\addplot[tangentred, very thick, domain=0.85:1.55]{1.2*(x - 1.2) + (0.5*1.2^2 + 0.5)};
\addplot[tangentred, very thick, domain=0.85:1.55]{1.2*(x - 1.2) + (0.5*1.2^2 + 1.5)};

\fill[tangentred] (axis cs:1.2, {0.5*1.2^2 - 1.5}) circle (1.8pt);
\fill[tangentred] (axis cs:1.2, {0.5*1.2^2 - 0.5}) circle (1.8pt);
\fill[tangentred] (axis cs:1.2, {0.5*1.2^2 + 0.5}) circle (2.2pt);
\fill[tangentred] (axis cs:1.2, {0.5*1.2^2 + 1.5}) circle (1.8pt);

\fill[tangentred] (axis cs:0, 0.5) circle (2.4pt)
  node[below right=3pt, curveblue, font=\scriptsize\bfseries] {$(0, 0.5)$};
\node[curveblue, font=\scriptsize\bfseries, anchor=center] at (axis cs:-1.5, 2.1) {$y = \tfrac{1}{2}t^2 + 0.5$};
\node[gray!70!black, font=\scriptsize, anchor=center] at (axis cs:-1.5, -1.4) {$y = \tfrac{1}{2}t^2 + C$};

\node[tangentred!90!black, font=\scriptsize\bfseries, align=center, anchor=north] at (axis cs:1.9, -1.1) {Parallel slopes\\$m = 1.2$};
\end{axis}
\end{tikzpicture}
\caption{The vertical freedom in a differential equation. For $\dv{y}{t}=t$, every solution is a parabola $y=\tfrac{1}{2}t^2 + C$. Along any vertical slice (such as $t=1.2$, purple dashed line), every parabola has the exact same slope ($m=1.2$, red tangents), because the rate depends only on $t$. Sliding a curve vertically leaves its slopes unchanged. An initial condition, like $(0,0.5)$, acts as a peg: it eliminates the vertical freedom and selects a single particular solution (blue).}
\label{ch36:fig-family-parabolas}
\end{figure}

This is the cleanest way to see why a single law yields a family (Figure~\ref{ch36:fig-family-parabolas}). The law fixes the \emph{shape} of motion -- the slopes everywhere -- but not the starting point. To select one solution from the family, we need exactly one more piece of data: an \emph{initial condition}.

\section{Initial conditions: pinning down one curve}
\label{ch36:sec-ic}

An \emph{initial condition} is a single anchoring fact of the form ``at this time, the quantity had this value,'' written
\[
  y(t_0) = y_0.
\]
It is the missing piece. The differential equation supplies the slopes; the initial condition supplies one point the curve must pass through. Together they typically determine exactly one solution. A differential equation packaged with an initial condition is called an \emph{initial value problem}, or IVP.

\begin{example}[An initial value problem]
\label{ch36:ex-ivp}
Solve
\[
  \dv{y}{t} = y, \qquad y(0) = 5.
\]
We already know the general solution is $y = Ce^t$. Now impose the initial condition: at $t=0$, we need $y = 5$. Substitute:
\[
  y(0) = Ce^{0} = C\cdot 1 = C = 5.
\]
So $C = 5$, and the unique solution to the initial value problem is
\[
  y = 5e^t.
\]
The family $Ce^t$ collapsed to a single curve the instant we said where it started. Geometrically, of all the exponential curves obeying the law, exactly one passes through the point $(0,5)$ -- and that is our answer.
\end{example}

\begin{focusbox}
The two-step rhythm of solving an initial value problem is worth committing to memory, because it recurs in every chapter to come:
\begin{enumerate}
  \item \textbf{Find the general solution} -- the whole family, with its free constant(s). This is the hard, creative step, the one that needs the techniques of later chapters.
  \item \textbf{Apply the initial condition(s)} to solve for the constant(s). This is the easy, mechanical step -- plug in the known point and do algebra.
\end{enumerate}
Mixing these up -- trying to satisfy the initial condition before you have the general family -- is the single most common source of confusion for beginners.
\end{focusbox}

\begin{trapbox}
A tempting error: thinking the initial condition has to be given ``at the beginning,'' at $t=0$. Not so. The condition $y(t_0)=y_0$ can anchor the curve at \emph{any} time $t_0$. ``Initial'' is a name, not a requirement. If a problem tells you $y(3) = 7$, you use that point exactly the same way -- substitute $t=3$, set the result equal to $7$, solve for the constant. The word ``initial'' is a historical leftover from physics, where the natural anchor is the state of a system when you start the clock.
\end{trapbox}

\begin{recoverybox}
\textbf{Pause and think:} If you are told that $\dv{y}{t} = y$ and $y(2) = 10$, what is the unique solution?
\emph{Mental check:} The general solution is $y(t) = Ce^t$. Setting $t=2$ yields $Ce^2 = 10$, so $C = 10e^{-2}$. The particular solution is therefore $y(t) = 10e^{-2}e^t = 10e^{t-2}$. Notice how cleanly the exponent shifts: when $t=2$, $e^{2-2} = 1$, instantly verifying $y(2) = 10$.
\end{recoverybox}

\subsection{Why we expect exactly one}

Pause on \emph{why} a single initial condition should select a single curve, because the picture here is so persuasive. Imagine you are a tiny walker dropped onto the plane at the starting point $(t_0, y_0)$. The differential equation is a compass: at your current location it tells you which direction to face -- it gives you a slope. You take an infinitesimal step in that direction. Now you are at a slightly new point; consult the compass again; it gives a new slope; step again. Marching forward (and backward) this way, you trace out a curve, and at every point of it your heading is exactly the slope the law prescribes -- so the curve is a solution. And the curve is \emph{forced}: at each instant you had no choice about which way to step. Start at the same point and you trace the same curve. \emph{One starting point, one path.}

That heuristic is essentially correct, and it is the seed of both a numerical method (Euler's method, which literally takes those small steps) and a theorem (the existence and uniqueness theorem of Section~\ref{ch36:sec-exist}). But the heuristic also quietly assumes the compass is well-behaved -- that the slope does not do anything pathological. We will see in a moment that when the law misbehaves, uniqueness can genuinely fail, and the picture of the forced walker breaks down in an instructive way.

\begin{figure}[H]
\centering
\begin{tikzpicture}[scale=0.92,>={Stealth[length=4pt]}]
  % Panel (a): The Local Compass
  \begin{scope}[shift={(0,0)}]
    \draw[fill=blue!3,draw=blue!15,rounded corners=4pt] (-2.6,-2.4) rectangle (2.6,2.75);
    \node[font=\small\bfseries,anchor=north] at (0,2.65) {(a) The Compass: Local Slopes};
    \draw[->,gray!50] (-2.2,-1.8)--(2.2,-1.8) node[right]{\scriptsize $t$};
    \draw[->,gray!50] (-2.0,-2.0)--(-2.0,1.7) node[above]{\scriptsize $y$};
    % grid of slope dashes
    \foreach \tx in {-1.4,-0.6,0.2,1.0,1.8}{
      \foreach \ty in {-1.2,-0.4,0.4,1.2}{
        \draw[gray!70,thick] (\tx-0.18,{\ty-0.18*(\ty*0.7+0.2)}) -- (\tx+0.18,{\ty+0.18*(\ty*0.7+0.2)});
      }
    }
    % Highlighted compass at one point
    \filldraw[fill=quotientorange!20,draw=quotientorange,thick] (0.2,0.4) circle (0.42cm);
    \draw[->,quotientorange!90!black,ultra thick] (0.2,0.4) -- ++(0.35,0.22);
    \node[quotientorange!90!black,font=\scriptsize\bfseries,anchor=south west] at (0.35,0.7) {slope $f(t,y)$};
    \node[font=\scriptsize,align=center] at (0,-2.1) {Local slope rule at every point};
  \end{scope}

  % Panel (b): The Walker Tracing the Unique Curve
  \begin{scope}[shift={(5.8,0)}]
    \draw[fill=blue!3,draw=blue!15,rounded corners=4pt] (-2.6,-2.4) rectangle (2.6,2.75);
    \node[font=\small\bfseries,anchor=north] at (0,2.65) {(b) The Walker: One Path};
    \draw[->,gray!50] (-2.2,-1.8)--(2.2,-1.8) node[right]{\scriptsize $t$};
    \draw[->,gray!50] (-2.0,-2.0)--(-2.0,1.7) node[above]{\scriptsize $y$};
    % background faint dashes
    \foreach \tx in {-1.4,-0.6,0.2,1.0,1.8}{
      \foreach \ty in {-1.2,-0.4,0.4,1.2}{
        \draw[gray!35,thin] (\tx-0.18,{\ty-0.18*(\ty*0.7+0.2)}) -- (\tx+0.18,{\ty+0.18*(\ty*0.7+0.2)});
      }
    }
    % Initial anchor point
    \coordinate (P0) at (-1.4,-0.5);
    \fill[tangentred] (P0) circle (2.2pt);
    \node[tangentred!90!black,font=\scriptsize\bfseries,below] at (P0) {$(t_0,y_0)$};
    % Stepping arrows
    \draw[->,tangentred,thick] (P0) -- ++(0.8,0.2) coordinate (P1);
    \fill[tangentred] (P1) circle (1.5pt);
    \draw[->,tangentred,thick] (P1) -- ++(0.8,0.35) coordinate (P2);
    \fill[tangentred] (P2) circle (1.5pt);
    \draw[->,tangentred,thick] (P2) -- ++(0.8,0.55) coordinate (P3);
    % Smooth solution curve
    \draw[curveblue,ultra thick] (-1.4,-0.5) .. controls (-0.6,-0.3) and (0.2,0.05) .. (1.8,1.4)
      node[right,font=\scriptsize\bfseries] {$y(t)$};
    \node[font=\scriptsize,align=center] at (0,-2.1) {One start $\implies$ one forced trajectory};
  \end{scope}
\end{tikzpicture}
\caption{The walker and the compass. (a) A differential equation $\dv{y}{t}=f(t,y)$ acts as
a local compass, assigning a prescribed heading to every coordinate in the plane.
(b) Dropping a walker at an initial point $(t_0,y_0)$ fixes the entire journey: stepping
along the compass directions forces out a single, unique solution curve $y(t)$.}
\label{ch36:fig-walker}
\end{figure}

\section{Order and linearity: classifying equations}
\label{ch36:sec-classify}

Differential equations take many forms. Two primary ways of classifying them are by \emph{order} and by \emph{linearity}. These two properties indicate the structure of an equation and the appropriate methods for solving it.

\begin{definition}[Order]
\label{ch36:def-order}
The \emph{order} of a differential equation is the order of the highest derivative that appears in it.
\end{definition}

Thus $\dv{y}{t} = y$ is \emph{first order} -- only the first derivative shows up. Newton's law of motion $m\,\frac{\dd^2 x}{\dd t^2} = F$, involving the second derivative (acceleration), is \emph{second order}. The order matters enormously: a first-order equation needs one initial condition to pin down a solution, a second-order equation needs \emph{two} (typically an initial position \emph{and} an initial velocity), and in general an $n$th-order equation carries $n$ free constants in its general solution and needs $n$ conditions to fix them.

\begin{connectionbox}
The rule ``order $n$ means $n$ free constants'' is the differential-equation echo of a fact you already know from integration. Each time you integrate, you collect one constant of integration. A second-order equation is, very loosely, ``two integrations deep,'' so it accumulates two constants; an $n$th-order equation, $n$ of them. This is exactly why pinning down a thrown ball's motion ($\frac{\dd^2 x}{\dd t^2} = -g$, second order) requires two facts -- where you threw it from and how fast -- not one. Integrating $-g$ once gives velocity with a constant; integrating again gives position with a second constant. Two integrations, two constants, two initial conditions.
\end{connectionbox}

The second classifier, linearity, is subtler but matters even more, because the entire theory splits along this line. Linear equations are the ones we can largely \emph{solve}; nonlinear ones are, in general, beyond closed-form reach and must be studied qualitatively or numerically.

\begin{definition}[Linear ODE]
\label{ch36:def-linear}
A differential equation is \emph{linear} if the unknown function $y$ and \emph{all} of its derivatives appear only to the first power, never multiplied together, and never inside any nonlinear function. A first-order linear ODE can always be written in the standard form
\[
  \dv{y}{t} + p(t)\,y = q(t),
\]
where $p$ and $q$ are given functions of $t$ alone. If the unknown appears in any other way -- squared, as $y^2$; multiplied by its own derivative, as $y\,\dv{y}{t}$; or inside a function, as $\sin y$ or $e^{y}$ -- the equation is \emph{nonlinear}.
\end{definition}

In a linear equation, the unknown and its derivatives appear only to the first power, scaled by coefficients that may depend on $t$ but not on $y$. Linearity is a condition on how the dependent variable $y$ enters, not on how $t$ enters.

\begin{example}[Sorting a handful of equations]
\label{ch36:ex-classify}
Classify a few equations to make the definitions concrete.
\begin{itemize}
  \item $\dv{y}{t} = y$ \quad -- first order, linear ($p(t) = -1$, $q(t)=0$).
  \item $\dv{y}{t} = t^2 y + \cos t$ \quad -- first order, linear; rewrite as $\dv{y}{t} - t^2 y = \cos t$, so $p = -t^2$, $q = \cos t$. The coefficient functions $p(t)$ and $q(t)$ may be nonlinear in $t$.
  \item $\dv{y}{t} = y^2$ \quad -- first order, \emph{nonlinear}: $y$ appears squared.
  \item $\dv{y}{t} = \sin y$ \quad -- first order, \emph{nonlinear}: $y$ sits inside $\sin$.
  \item $\frac{\dd^2 x}{\dd t^2} + 3\dv{x}{t} + 2x = 0$ \quad -- second order, linear.
  \item $\frac{\dd^2 \theta}{\dd t^2} + \sin\theta = 0$ \quad -- second order, \emph{nonlinear} (the pendulum equation; $\sin\theta$ is the culprit).
\end{itemize}
\end{example}

\begin{recoverybox}
\textbf{Pause and think:} Why is $\dv{y}{t} = t^4 y$ a linear equation, whereas $\dv{y}{t} = t y^4$ is nonlinear? The power $4$ appears in both!
\emph{Mental check:} Linearity cares strictly about the \emph{unknown function} $y$ and its derivatives, never about the independent variable $t$. In $\dv{y}{t} = t^4 y$, the unknown $y$ appears purely to the first power; the factor $t^4$ is merely a time-dependent coefficient $p(t) = -t^4$. In $\dv{y}{t} = t y^4$, the unknown itself is raised to the fourth power, which immediately breaks linearity.
\end{recoverybox}

\begin{aestheticbox}
The central role of linearity stems from the \emph{superposition principle}: if $y_1$ and $y_2$ both solve a linear \emph{homogeneous} equation (where $q(t)=0$), then so does every linear combination $c_1 y_1 + c_2 y_2$. This algebraic fact -- that the solution set forms a vector space -- is the foundation for the theory of linear ODEs and enables systematic solution formulas. Nonlinear equations have no such principle; their solutions cannot be combined by addition. Superposition makes linearity the primary dividing line of the subject.
\end{aestheticbox}

\section{The slope field: seeing solutions before solving}
\label{ch36:sec-slopefield}

Here is the idea that turns an abstract equation into a picture you can read with your eyes. It rests entirely on the mechanism exposed earlier: a first-order equation $\dv{y}{t} = f(t,y)$ assigns a \emph{slope to every point} of the plane. So \emph{draw} those slopes.

At a grid of points $(t,y)$, compute the number $f(t,y)$ and draw a short line segment through the point with that slope. Do this at many points and you fill the plane with a texture of tiny tilted dashes -- a \emph{slope field} (also called a direction field). Each dash is a local instruction: ``a solution passing through here must be heading this way.'' A solution curve is then any curve that flows along the dashes, tangent to the field at every point it touches, like a leaf carried by a current.

\begin{connectionbox}
\textbf{The geometric turn of Poincar\'e.}
The idea of reading an equation through its field of directions is older than it looks. The eighteenth-century analysts already thought of $\dv{y}{t}=f(t,y)$ as prescribing a direction at each point, but it was the qualitative theory of differential equations -- pioneered by Henri Poincar\'e in the 1880s in his memoirs on curves defined by differential equations -- that made the geometric viewpoint a discipline in its own right. Faced with equations that no algebraic formula could ever solve, Poincar\'e shifted the question from ``can we write down an exact formula for the solution?'' to ``what does the family of solutions \emph{look like}?'' -- where the curves go, which ones are stable, and what happens in the long run. The slope field is the gateway to that whole way of thinking, and it is the seed of the phase-plane analysis you will meet in Chapter~39.
\end{connectionbox}

\begin{intuitionbox}
The slope field is the differential equation made \emph{visible}. You never solve anything; you just plot the law's slope at a mesh of points. And yet, once the field is drawn, the solutions reveal themselves to the eye: trace any curve that stays tangent to the little dashes, and you are tracing a solution. Different starting heights give different curves -- the family of solutions appears as a family of streamlines combing through the field.
\end{intuitionbox}

Take the slope field for our very first equation, $\dv{y}{t} = y$. Here $f(t,y) = y$: the prescribed slope at a point is simply its height. Above the $t$-axis (where $y>0$) slopes are positive and steepen as you rise; below it (where $y<0$) slopes are negative and steepen as you fall; right on the axis ($y=0$) the slope is $0$, flat. Figure~\ref{ch36:fig-slopefield-exp} shows the result, with three solution curves $y = Ce^t$ threaded through it.

\begin{figure}[H]
\centering
\begin{tikzpicture}
\begin{axis}[
  width=12cm, height=8cm,
  xlabel={$t$}, ylabel={$y$},
  xlabel style={anchor=south west},
  xmin=-2.2, xmax=2.2, ymin=-3.2, ymax=3.2,
  axis lines=middle,
  xtick={-2,1,2}, ytick={-3,-2,2,3},
  xticklabel style={fill=white, inner sep=1pt, font=\footnotesize},
  yticklabel style={fill=white, inner sep=1pt, font=\footnotesize},
  axis on top,
  clip mode=individual,
  enlargelimits=false,
  clip=true,
]
% slope field: at (t,y) slope = y. Draw normalized segments.
\foreach \t in {-2,-1.5,-1,-0.5,0,0.5,1,1.5,2}{
  \foreach \yy in {-3,-2.5,-2,-1.5,-1,-0.5,0,0.5,1,1.5,2,2.5,3}{
    \edef\sl{\yy}
    \addplot[gray, thick, samples=2, domain=\t-0.16:\t+0.16, forget plot]
      ({x}, {\yy + (\sl)*(x-\t)});
  }
}
% three solution curves y = C e^t
\addplot[curveblue, very thick, samples=120, domain=-2.2:1.18]{exp(x)};
\node[curveblue, font=\scriptsize\bfseries, fill=white, inner sep=1.5pt, rounded corners=1pt, anchor=south east] at (axis cs:0.5, 2.1) {$y=e^t$};
\addplot[tangentred, very thick, samples=120, domain=-2.2:2.05]{0.4*exp(x)};
\node[tangentred, font=\scriptsize\bfseries, fill=white, inner sep=1.5pt, rounded corners=1pt, anchor=north east] at (axis cs:2.0, 2.2) {$y=0.4e^t$};
\addplot[targetgreen, very thick, samples=120, domain=-2.2:1.18]{-exp(x)};
\node[targetgreen, font=\scriptsize\bfseries, fill=white, inner sep=1.5pt, rounded corners=1pt, anchor=north east] at (axis cs:0.5, -2.1) {$y=-e^t$};
\node[gray!70!black, font=\scriptsize\itshape, fill=white, inner sep=1pt, anchor=south west] at (axis cs:-2.1,0.15) {$y=0$ (flat)};
\end{axis}
\end{tikzpicture}
\caption{The slope field for $\dv{y}{t}=y$. Each gray dash has slope equal to its own height $y$. Three solution curves $y=Ce^t$ flow tangent to the field: $C=1$ (blue), $C=0.4$ (red), and $C=-1$ (green). Notice how every curve hugs the dashes everywhere it passes.}
\label{ch36:fig-slopefield-exp}
\end{figure}

Spend a moment really looking at Figure~\ref{ch36:fig-slopefield-exp}. The dashes are flat along the $t$-axis and tilt ever more steeply as you move away from it -- exactly because the slope equals the height. The blue curve, $y=e^t$, rises through the steepening forest of positive slopes, always tangent. The green curve, $y=-e^t$, plunges through the negative slopes below. And the $t$-axis itself, $y=0$, is a perfectly flat solution -- the constant zero function, with slope $0$ everywhere, matching the flat dashes that lie along it. You are reading the entire solution family off a picture, having solved nothing.

\begin{walkbox}
Look again at Figure~\ref{ch36:fig-slopefield-exp} and ask yourself: through the single point $(0, 0.4)$, how many solution curves pass? Trace the dashes by eye starting there. You should find that the curve is forced -- there is exactly one streamline through that point (the red one, $y=0.4e^t$), because at every step the field dictates your heading. This is the geometric face of uniqueness: a well-behaved slope field carries exactly one solution curve through each point. Now ask the harder question: can two \emph{different} solution curves ever cross each other? If they did, they would share a point -- and at that shared point the field prescribes a single slope, so both curves would have to leave it in the same direction. Solution curves of a nice first-order equation never cross. That is uniqueness, seen.
\end{walkbox}

\subsection{Reading a field with a \texorpdfstring{$t$}{t}-dependent slope}

The previous field had slopes depending only on $y$, so the pattern looked the same at every horizontal position. Now take one where the clock matters too:
\[
  \dv{y}{t} = t - y.
\]
Now the prescribed slope at $(t,y)$ is $t - y$, which changes as you move both horizontally and vertically. Where $t = y$ -- along the diagonal line $y=t$ -- the slope is zero, so the dashes lie flat there. Above that diagonal ($y > t$) slopes are negative; below it ($y<t$) slopes are positive. Everything is gently herded toward a special straight solution. Figure~\ref{ch36:fig-slopefield-ty} shows the field with two solution curves.

\begin{figure}[H]
\centering
\begin{tikzpicture}
\begin{axis}[
  width=12cm, height=8cm,
  xlabel={$t$}, ylabel={$y$},
  xlabel style={anchor=south west},
  xmin=-3.2, xmax=3.2, ymin=-3.2, ymax=3.2,
  axis lines=middle,
  xtick={-2,-1,1,2}, ytick={-2,-1,1,2},
  xticklabel style={fill=white, inner sep=1pt, font=\footnotesize},
  yticklabel style={fill=white, inner sep=1pt, font=\footnotesize},
  axis on top,
  clip mode=individual,
  enlargelimits=false,
  clip=true,
]
\foreach \t in {-2.5,-2,-1.5,-1,-0.5,0,0.5,1,1.5,2,2.5}{
  \foreach \yy in {-2.5,-2,-1.5,-1,-0.5,0,0.5,1,1.5,2,2.5}{
    \edef\sl{\t-\yy}
    \addplot[gray, thick, samples=2, domain=\t-0.14:\t+0.14, forget plot]
      ({x}, {\yy + (\sl)*(x-\t)});
  }
}
% the particular straight solution y = t - 1
\addplot[structurepurple, very thick, dashed, samples=2, domain=-2.2:3.2]{x-1};
\node[structurepurple, font=\scriptsize\bfseries, fill=white, inner sep=1.5pt, rounded corners=1pt, anchor=south east] at (axis cs:-0.2,-0.9) {attractor $y=t-1$};
% two solution curves of y' = t - y : general y = t - 1 + C e^{-t}
\addplot[curveblue, very thick, samples=140, domain=-1.4:3.2]{x-1+2.5*exp(-x)};
\node[curveblue, font=\scriptsize\bfseries, fill=white, inner sep=1.5pt, rounded corners=1pt, anchor=south east] at (axis cs:2.6,2.1) {$y = t - 1 + 2.5e^{-t}$};
\addplot[tangentred, very thick, samples=140, domain=-1.0:3.2]{x-1-2.5*exp(-x)};
\node[tangentred, font=\scriptsize\bfseries, fill=white, inner sep=1.5pt, rounded corners=1pt, anchor=north west] at (axis cs:1.8,-0.7) {$y = t - 1 - 2.5e^{-t}$};
\end{axis}
\end{tikzpicture}
\caption{The slope field for $\dv{y}{t}=t-y$. Dashes are flat along the diagonal $y=t$ and tilt away from it. The purple dashed line $y=t-1$ is itself a solution -- an \emph{attractor} every other solution curls toward. The blue and red curves, $y=t-1\pm 2.5e^{-t}$, both rush in to merge with it as $t$ grows.}
\label{ch36:fig-slopefield-ty}
\end{figure}

There is a striking feature in Figure~\ref{ch36:fig-slopefield-ty} that you could not have guessed from the bare equation but that the picture announces immediately: the line $y = t-1$ is itself a solution, and it acts as an \emph{attractor}. Check that it solves the equation: if $y = t-1$ then $\dv{y}{t} = 1$, while the right side is $t - y = t - (t-1) = 1$. They match. And every other solution, no matter where it starts, bends toward this line as $t$ increases -- you can watch the blue and red curves swoop in and shadow it. This long-run behavior is exactly the kind of qualitative insight slope fields hand you for free, long before you can write the general solution $y = t - 1 + Ce^{-t}$ (which, you will be able to confirm, this is).

\begin{happeningbox}
What is really happening when you ``read a slope field'' is that you are performing, by eye, the same construction the existence theorem performs by logic and Euler's method performs by arithmetic. All three reconstruct a global solution from purely local slope data. Your eye integrates the field into a smooth streamline; Euler's method integrates it into a polygonal approximation by taking literal small steps; the theorem proves that this integration always succeeds (under mild hypotheses) and yields exactly one curve per starting point. The slope field is therefore not a mere illustration -- it is the geometric content of the equation itself. An equation you cannot solve in formulas, you can very often \emph{understand completely} from its field.
\end{happeningbox}

\subsection{A field we cannot solve, but can read: the logistic law}
\label{ch36:sub-logistic-field}

The two fields so far belonged to equations we could solve in closed form, so the slope field merely confirmed what algebra already knew. The geometric perspective is especially useful when algebraic solutions are difficult to obtain. Here is the cleanest example in elementary modeling -- the \emph{logistic equation}, the corrected law of population growth that we will meet again in Section~\ref{ch36:sec-model}:
\[
  \dv{P}{t} = k\,P\!\left(1 - \frac{P}{M}\right).
\]
The right-hand side is nonlinear (notice the $P^2$ term), so standard linear formulas do not apply, and an explicit solution requires separation of variables. But its slope field can be drawn directly, displaying the qualitative behavior.

Read the law as a rule for slope. Two heights make the right side vanish: $P=0$ and $P=M$. At those two heights the prescribed slope is zero, so the dashes lie perfectly flat -- they mark two \emph{equilibrium} (constant) solutions, two horizontal lines where the slope is zero. Between them, for $0 < P < M$, both factors $P$ and $1-P/M$ are positive, so the slope is positive and every solution climbs. Above the carrying capacity, $P > M$, the factor $1-P/M$ turns negative, the slope turns negative, and solutions fall. Both regions drain toward the line $P=M$: it is an attractor, just like the line $y=t-1$ was for the earlier field. Figure~\ref{ch36:fig-logistic} draws it, taking $k=1$ and $M=3$.

\begin{figure}[H]
\centering
\begin{tikzpicture}
\begin{axis}[
  width=12cm, height=8cm,
  xlabel={$t$}, ylabel={$P$},
  xlabel style={anchor=south west},
  ylabel style={anchor=south west},
  xmin=0, xmax=6, ymin=-0.4, ymax=4.4,
  axis lines=middle,
  xtick={1,2,3,4,5,6}, ytick={1,2,3,4},
  enlargelimits=false,
  clip=true,
]
% slope field for P' = P(1 - P/3), k=1, M=3
\foreach \t in {0.5,1.1,1.7,2.3,2.9,3.5,4.1,4.7,5.3,5.9}{
  \foreach \pp in {0,0.4,0.8,1.2,1.6,2,2.4,2.8,3,3.2,3.6,4}{
    \edef\sl{\pp*(1-\pp/3)}
    \addplot[gray, thick, samples=2, domain=\t-0.18:\t+0.18, forget plot]
      ({x}, {\pp + (\sl)*(x-\t)});
  }
}
% equilibrium lines P=0 and P=M=3
\addplot[structurepurple, very thick, dashed, samples=2, domain=0:6]{3};
\node[structurepurple, font=\scriptsize\bfseries, fill=white, inner sep=1.5pt, rounded corners=1pt, anchor=north east] at (axis cs:3.8, 2.7) {carrying capacity $M=3$};
\addplot[structurepurple, very thick, dashed, samples=2, domain=0:6]{0};
\node[structurepurple, font=\scriptsize\bfseries, fill=white, inner sep=1.5pt, rounded corners=1pt, anchor=south east] at (axis cs:5.5, 0.12) {extinction $P=0$};
% logistic solution from below: P0 = 0.3.  P(t) = 3/(1 + 9 e^{-t})
\addplot[curveblue, very thick, samples=160, domain=0:6]{3/(1 + 9*exp(-x))};
\node[curveblue, font=\scriptsize\bfseries, fill=white, inner sep=1.5pt, rounded corners=1pt, anchor=north west] at (axis cs:2.7, 1.4) {$P_0 = 0.3$};
% solution coming down from above: P0 = 4.  P(t) = 3/(1 - 0.25 e^{-t})
\addplot[tangentred, very thick, samples=160, domain=0:6]{3/(1 - 0.25*exp(-x))};
\node[tangentred, font=\scriptsize\bfseries, fill=white, inner sep=1.5pt, rounded corners=1pt, anchor=south west] at (axis cs:1.2, 3.5) {$P_0 = 4$};
\end{axis}
\end{tikzpicture}
\caption{The slope field for the logistic law $\dv{P}{t}=P(1-P/3)$ (so $k=1$, $M=3$). The dashes lie flat along the two equilibria $P=0$ and $P=M=3$ (purple dashed). The blue curve starts small ($P_0=0.3$) and traces the famous ``S'' shape -- slow, then fast, then leveling off at the carrying capacity. The red curve starts above capacity ($P_0=4$) and sinks toward the same line. Every solution is herded toward $P=M$.}
\label{ch36:fig-logistic}
\end{figure}

\begin{walkbox}
Follow the blue curve in Figure~\ref{ch36:fig-logistic} across three stages of growth. Near the bottom the dashes are nearly flat (small $P$ makes the slope small), so the curve creeps -- a young population grows slowly because there are few to reproduce. Through the middle, around $P=M/2 = 1.5$, the slope is steepest (the product $P(1-P/M)$ is maximal there), and the curve shoots up -- this is the rapid growth phase. Near the top the dashes flatten again as $P$ nears $M$, and the curve eases into a plateau -- growth slows as the population approaches capacity. Slow, fast, slow: that produces the characteristic sigmoid curve. We have determined the qualitative behavior of a self-limiting population without solving the nonlinear equation analytically. The red curve from above similarly approaches the attractor $P=M$ from higher initial values.
\end{walkbox}

This illustrates the utility of slope fields. The logistic equation is a nonlinear law that can be solved analytically by separation of variables (as shown in the next chapter) -- but even before doing so, the direction field displays the equilibria, their stability, and the characteristic shape of every solution. For equations that resist closed-form solution entirely, this geometric approach is essential.

\section{Existence and uniqueness, intuitively}
\label{ch36:sec-exist}

We have been assuming, on the strength of the walker picture, that an initial value problem has exactly one solution. It is time to say carefully when that is true -- and, just as instructively, when it fails. The governing result is one of the cornerstones of the subject.

\begin{theorem}[Existence and Uniqueness, informal]
\label{ch36:thm-exist}
Consider the initial value problem
\[
  \dv{y}{t} = f(t,y), \qquad y(t_0) = y_0.
\]
If $f$ is continuous in a rectangle around the point $(t_0, y_0)$, then at least one solution exists through that point (\emph{existence}). If, in addition, the partial derivative $\pdv{f}{y}$ is continuous there, then the solution is \emph{unique} -- exactly one solution curve passes through $(t_0,y_0)$.
\end{theorem}

The theorem has two halves, and they require different amounts of niceness from the law $f$. Mere continuity of $f$ buys you \emph{existence}: some solution gets going. To also guarantee \emph{uniqueness} -- that the walker's path is forced, with no branching -- you need the extra smoothness that $\pdv{f}{y}$ be continuous (a slightly weaker ``Lipschitz'' condition is what the full proof really uses, but continuity of $\pdv{f}{y}$ is the clean sufficient version). Intuitively, $\pdv{f}{y}$ controls how fast the prescribed slope changes as you move up and down; keeping it under control prevents nearby curves from being able to split apart or merge.

\begin{intuitionbox}
Read the theorem as a guarantee about the walker. \emph{Existence} says: if the compass (the law $f$) points somewhere sensible at and near your starting spot, you can always start walking -- a path exists. \emph{Uniqueness} says: if the compass also turns \emph{smoothly} as you move (no sudden infinite swings in the prescribed slope), then your path is completely determined -- there is no fork in the road, no moment where the law leaves you a choice. Together they justify the whole enterprise: a sensible law plus a starting point yields one and only one future.
\end{intuitionbox}

\begin{recoverybox}
\textbf{Pause and think:} Does Theorem~\ref{ch36:thm-exist} guarantee that a unique solution passes through $(0, 2)$ for the equation $\dv{y}{t} = 3y + \cos t$?
\emph{Mental check:} Here $f(t,y) = 3y + \cos t$, which is continuous everywhere. Its partial derivative with respect to $y$ is $\pdv{f}{y} = 3$, which is a constant and therefore continuous across the entire plane. Both hypotheses hold everywhere, so exactly one solution curve passes through $(0,2)$---and through any other point you choose!
\end{recoverybox}

Without this theorem, deterministic modeling lacks a formal guarantee. If a physical process is modeled by a differential equation and an initial condition, we expect exactly one determined outcome. The existence and uniqueness theorem confirms this mathematically, ensuring that a well-posed initial value problem yields exactly one solution.

\subsection{When uniqueness fails -- and why that is not a paradox}

The following example illustrates why the hypotheses of Theorem~\ref{ch36:thm-exist} are necessary. Consider the problem
\[
  \dv{y}{t} = \sqrt{\abs{y}}, \qquad y(0) = 0.
\]
Here $f(t,y) = \sqrt{\abs{y}}$ is continuous everywhere, so existence is satisfied -- a solution exists. But watch what $\pdv{f}{y}$ does near $y=0$. For $y>0$, $f = \sqrt{y}$ and $\pdv{f}{y} = \tfrac{1}{2\sqrt{y}}$, which blows up to infinity as $y\to 0$. The uniqueness hypothesis fails at the starting point, and uniqueness itself fails with it:

\begin{example}[Two solutions from one start]
\label{ch36:ex-nonunique}
The constant function $y(t) = 0$ solves the problem: its derivative is $0$, and $\sqrt{\abs{0}} = 0$, so the equation holds, and $y(0)=0$. But the function
\[
  y(t) = \tfrac{1}{4}t^2 \quad (t \ge 0)
\]
\emph{also} solves it. Check: $\dv{y}{t} = \tfrac{1}{2}t$, while $\sqrt{\abs{y}} = \sqrt{\tfrac{1}{4}t^2} = \tfrac{1}{2}t$ for $t\ge 0$. The two sides match, and $y(0)=0$. So \emph{two different} solution curves emerge from the single point $(0,0)$ -- the flat one that stays put, and the rising parabola. In fact there are infinitely many: the curve can linger at zero for any length of time $a\ge 0$ and then peel up into the shifted parabola $y=\tfrac14(t-a)^2$. Figure~\ref{ch36:fig-branch} shows a few of them fanning out of the origin. Uniqueness does not hold here.
\end{example}

\begin{figure}[H]
\centering
\begin{tikzpicture}
\begin{axis}[
  width=11cm, height=6cm,
  xlabel={$t$}, ylabel={$y$},
  xlabel style={anchor=south west},
  xmin=-0.4, xmax=4.2, ymin=-0.4, ymax=4.2,
  axis lines=middle,
  xtick={1,2,3,4}, ytick={1,2,3,4},
  enlargelimits=false,
  clip=true,
]
% the flat solution y = 0
\addplot[targetgreen, very thick, samples=2, domain=0:4]{0};
\node[targetgreen, font=\scriptsize\bfseries, anchor=south west] at (axis cs:0.6,0.15) {$y=0$};
% parabolas that linger at 0 until t=a then rise: y = (1/4)(t-a)^2
\addplot[curveblue, very thick, samples=80, domain=0:4]{0.25*x^2};
\node[curveblue, font=\scriptsize\bfseries, anchor=south east] at (axis cs:2.1,1.8) {$y=\tfrac{1}{4}t^2$};
\addplot[tangentred, very thick, samples=70, domain=1:4]{0.25*(x-1)^2};
\node[tangentred, font=\scriptsize\bfseries, anchor=north west] at (axis cs:3.0,0.85) {$a=1$};
\addplot[structurepurple, very thick, samples=60, domain=2:4]{0.25*(x-2)^2};
\node[structurepurple, font=\scriptsize\bfseries, anchor=north west] at (axis cs:3.6,0.4) {$a=2$};
\fill[tangentred] (axis cs:0,0) circle (2.2pt);
\node[tangentred!90!black, font=\scriptsize\bfseries, anchor=north east] at (axis cs:-0.05,-0.05) {$(0,0)$};
\end{axis}
\end{tikzpicture}
\caption{Non-uniqueness for $\dv{y}{t}=\sqrt{\abs{y}}$, $y(0)=0$. Infinitely many solution curves leave the single starting point $(0,0)$: the flat solution $y=0$ (green), and a whole fan of parabolas $y=\tfrac14(t-a)^2$ that linger at zero until an arbitrary time $a$ and then rise. The uniqueness hypothesis $\pdv{f}{y}$ continuous fails at $y=0$, and this is what failure looks like.}
\label{ch36:fig-branch}
\end{figure}

\begin{trapbox}
The lesson is that the hypotheses of the theorem are essential.
When $\pdv{f}{y}$ is continuous, uniqueness is guaranteed; when that condition fails (here, an infinite partial derivative at $y=0$), solutions can branch. Never assume uniqueness for a nonlinear equation without checking the hypothesis. The square-root law above is not an artificial curiosity: it is essentially the equation governing how fast a leaking tank drains (Torricelli's law), and the non-uniqueness has a physical reading -- you cannot run the drainage backward in time to recover when the tank was full, because many different fill-times lead to the same empty tank.
\end{trapbox}

\section{Why differential equations exist: modeling}
\label{ch36:sec-model}

We have treated differential equations as mathematical objects, but they earn their central place because they are the language of \emph{change in the physical world}. A modeling problem almost always arrives as a statement about rates, and translating that statement into symbols produces a differential equation. We construct three classical models from basic physical principles.

\subsection{Population growth}

\begin{applicationbox}
\textbf{The sentence:} \emph{A population grows at a rate proportional to its current size.} The reasoning is straightforward -- twice as many individuals means twice as many births per year. Let $P(t)$ be the population. ``Grows at a rate'' is $\dv{P}{t}$; ``proportional to its current size'' means that rate is some constant $k$ times $P$. So
\[
  \dv{P}{t} = kP.
\]
This is the differential equation $\dv{P}{t}=kP$ applied to population dynamics. Its solution is $P(t) = P_0 e^{kt}$, where $P_0 = P(0)$ is the starting population: \emph{exponential growth}, the model that explains why unchecked populations, compound interest, and chain reactions all grow exponentially. The constant $k$ is the per-capita growth rate; a larger $k$ means faster growth.
\end{applicationbox}

Exponential growth is, of course, too optimistic forever -- no real population grows without bound. The fix is itself a lesson in modeling: add a term that brakes the growth as $P$ approaches a carrying capacity $M$. That gives the \emph{logistic equation} $\dv{P}{t} = kP(1 - P/M)$, a nonlinear model whose slope field shows every solution leveling off at $M$. We will not solve it here, but notice how naturally a refined sentence (``growth slows as the population fills its niche'') produces a refined equation.

\subsection{Newton's law of cooling}

\begin{applicationbox}
\textbf{The sentence:} \emph{An object cools at a rate proportional to the gap between its temperature and the surrounding temperature.} Let $T(t)$ be the object's temperature and $A$ the ambient temperature. The ``gap'' is $T - A$. The rate of cooling is $\dv{T}{t}$, and ``proportional to the gap'' (with cooling, so the constant is negative) gives
\[
  \dv{T}{t} = -k\,(T - A), \qquad k > 0.
\]
When $T > A$ the right side is negative and the object cools; when $T < A$ it is positive and the object warms; when $T = A$ the rate is zero and the temperature holds. The solution is $T(t) = A + (T_0 - A)e^{-kt}$ -- the temperature decays exponentially toward the ambient $A$, exactly the attractor behavior we saw in Figure~\ref{ch36:fig-slopefield-ty}. This is the equation forensic scientists use to estimate time of death from body temperature.
\end{applicationbox}

\subsection{Motion under a force}

\begin{applicationbox}
\textbf{The sentence:} \emph{Force equals mass times acceleration.} Acceleration is the second derivative of position, $\frac{\dd^2 x}{\dd t^2}$, so Newton's second law is already a second-order differential equation:
\[
  m\,\frac{\dd^2 x}{\dd t^2} = F.
\]
For a stone in free fall the force is constant, $F = -mg$, and the equation collapses to $\frac{\dd^2 x}{\dd t^2} = -g$. Integrate once for velocity, again for position, collecting two constants -- which the two initial conditions (starting height and starting velocity) fix. For a mass on a spring the force is $F = -kx$ (Hooke's law, pulling back toward rest), giving $m\,\frac{\dd^2 x}{\dd t^2} = -kx$, whose solutions oscillate sinusoidally forever. The same skeleton, $F = ma$, with a different force law, produces falling, orbiting, oscillating, and swinging -- the entire repertoire of classical motion.
\end{applicationbox}

\begin{aestheticbox}
Notice the unity here. Three completely different phenomena -- a breeding population, a cooling cup, a falling stone -- led to differential equations, and two of them led to the \emph{same} equation, $\dv{y}{t} = ky$, up to the sign and meaning of the constant. This is what gives the subject its reach: a handful of equation-shapes recur across all of science, so that mastering them addresses problems across biology, physics, chemistry, and finance. The differential equation is the level of abstraction at which a population, a bank account, and a radioactive sample share the same mathematical structure.
\end{aestheticbox}

\begin{labbox}
Before reading on, try translating each sentence into a differential equation yourself -- do not solve, just write the equation.
\begin{enumerate}
  \item A radioactive sample decays at a rate proportional to the amount remaining.
  \item A skydiver's downward speed $v$ increases due to gravity but is opposed by air resistance proportional to the speed.
  \item A savings account grows by continuous interest at rate $r$, and money is also withdrawn at a steady $w$ dollars per year.
\end{enumerate}
Answers: (1) $\dv{N}{t} = -kN$; (2) $m\,\dv{v}{t} = mg - bv$; (3) $\dv{B}{t} = rB - w$. Notice that each is just a sentence about rates, set equal to symbols. That translation -- prose to equation -- is the real modeling skill, and it is almost entirely a matter of identifying ``rate of change of \underline{\hspace{1.5em}}'' with a derivative.
\end{labbox}

\section{Looking ahead}
\label{ch36:sec-ahead}

You now know what a differential equation is, what it means to solve one, how an initial condition selects a single solution from a family, how to classify equations by order and linearity, how to \emph{see} solutions in a slope field, when a solution is guaranteed to exist and be unique, and why these equations are the language in which the sciences write their laws. What remains is technique -- the actual machinery for producing the general solutions we have so far only verified or read off pictures.

The next three chapters supply that machinery. Chapter~37, \emph{First-Order Differential Equations}, immediately develops systematic methods for first-order equations -- separation of variables, which solves $\dv{y}{t}=ky$ and the logistic equation, and the integrating factor, which solves first-order linear equations $\dv{y}{t}+p(t)y=q(t)$. The chapter after that tackles second-order linear equations, where superposition gives a complete and elegant theory -- this is the home of springs, circuits, and resonance. And the final chapter pushes onward to systems of equations and the qualitative, geometric viewpoint, where the slope-field thinking of this chapter blossoms into the full phase-plane analysis of how coupled quantities evolve together. The pictures you have learned to read here are the ones you will be reading for the rest of the subject.

\section*{Exercises}
\addcontentsline{toc}{section}{Exercises}

\begin{exercise}[Mental check: equations and their unknowns]
\label{ch36:ex-mental1}
For each equation, identify what kind of mathematical object a solution must be (a number, a coordinate point, or an entire function):
(a) $x^2 - 4x + 3 = 0$; \quad
(b) $\dv{y}{t} = 3y - t$; \quad
(c) $x^2 + y^2 = 25$; \quad
(d) $\frac{\dd^2 y}{\dd t^2} + 9y = 0$.
Which of these are ordinary differential equations?
\end{exercise}

\begin{exercise}[Mental check: reading slope from the law]
\label{ch36:ex-mental2}
Without solving or finding an antiderivative, determine the slope of the solution curve to $\dv{y}{t} = 2t - y$ as it passes through the point $(t, y) = (1, 3)$. Is the solution curve rising, falling, or flat at that exact instant? Where in the $(t,y)$ plane would solution curves have horizontal tangents?
\end{exercise}

\begin{exercise}
\label{ch36:ex-verify1}
Verify by substitution that $y = \sin(2t)$ is a solution of the second-order equation $\frac{\dd^2 y}{\dd t^2} + 4y = 0$. Then verify that $y = \cos(2t)$ is also a solution, and that $y = 3\sin(2t) - 5\cos(2t)$ is too. What property of the equation makes this last combination work?
\end{exercise}

\begin{exercise}
\label{ch36:ex-classify1}
For each equation state its order, and say whether it is linear or nonlinear (give the reason for any nonlinearity):
(a) $\dv{y}{t} = 3y - 7$; \quad
(b) $\dv{y}{t} = y\cos t$; \quad
(c) $\dv{y}{t} = \cos(y)$; \quad
(d) $\frac{\dd^2 x}{\dd t^2} + x\,\dv{x}{t} + x = 0$; \quad
(e) $t^2\frac{\dd^2 y}{\dd t^2} - t\dv{y}{t} + y = e^t$.
\end{exercise}

\begin{exercise}
\label{ch36:ex-ivp1}
The general solution of $\dv{y}{t} = -2y$ is $y = Ce^{-2t}$. Find the particular solution satisfying $y(0) = 7$. Then find the one satisfying $y(1) = 7$ instead, and explain why the constant $C$ comes out differently in the two cases.
\end{exercise}

\begin{exercise}
\label{ch36:ex-integrate1}
Solve the initial value problem $\dv{y}{t} = 6t^2 - 4$, $y(0) = 3$, by direct integration. (This equation's right-hand side depends only on $t$, so a single antiderivative does the job.)
\end{exercise}

\begin{exercise}
\label{ch36:ex-slopefield1}
Consider $\dv{y}{t} = -y/t$ for $t > 0$.
(a) Along which curves in the plane is the prescribed slope equal to zero? Where is it undefined?
(b) Verify that $y = C/t$ is a solution for every constant $C$.
(c) Sketch, by hand, the slope field at a few points and confirm that the curves $y = C/t$ (hyperbolas) flow tangent to it.
\end{exercise}

\begin{exercise}
\label{ch36:ex-attractor1}
For Newton's law of cooling $\dv{T}{t} = -k(T - A)$ with $k>0$:
(a) What constant function $T(t)$ is a solution? (This is the \emph{equilibrium} temperature.)
(b) Using the slope-field reasoning of Section~\ref{ch36:sec-slopefield}, argue without solving that every solution approaches $A$ as $t\to\infty$.
(c) Verify that $T(t) = A + (T_0 - A)e^{-kt}$ solves the equation and satisfies $T(0) = T_0$.
\end{exercise}

\begin{exercise}
\label{ch36:ex-order1}
Explain, in your own words, why the general solution of an $n$th-order differential equation should contain $n$ arbitrary constants. Then state how many initial conditions are needed to determine a unique solution of (a) $\dv{y}{t} = y^2$, and (b) $m\frac{\dd^2 x}{\dd t^2} = -kx$, and say physically what those conditions represent in case (b).
\end{exercise}

\begin{exercise}
\label{ch36:ex-nonunique1}
Consider $\dv{y}{t} = 3y^{2/3}$ with $y(0) = 0$.
(a) Verify that $y(t) = 0$ is a solution.
(b) Verify that $y(t) = t^3$ is also a solution.
(c) Compute $\pdv{f}{y}$ for $f(t,y) = 3y^{2/3}$ and explain, by reference to Theorem~\ref{ch36:thm-exist}, why the existence of two solutions does not contradict the uniqueness theorem.
\end{exercise}

\begin{exercise}
\label{ch36:ex-model1}
Translate each sentence into a differential equation (do not solve):
(a) A tank holds $V$ litres of water; salt water flows in carrying salt at a constant $r$ grams per minute, and well-mixed water flows out carrying salt at a rate proportional to the current amount $S$ of salt.
(b) The velocity $v$ of a parachutist changes under gravity minus an air-resistance force proportional to $v^2$.
(c) A rumour spreads through a town of $N$ people at a rate proportional to the product of those who know it, $y$, and those who do not, $N - y$.
\end{exercise}

\begin{exercise}
\label{ch36:ex-logistic1}
The logistic equation is $\dv{P}{t} = kP\!\left(1 - \dfrac{P}{M}\right)$ with $k, M > 0$.
(a) Find the two constant (equilibrium) solutions by setting the right-hand side to zero.
(b) For $0 < P < M$, is $\dv{P}{t}$ positive or negative? What does this say about a population that starts below carrying capacity?
(c) Without solving, sketch the qualitative shape of a solution that starts at $P(0)$ slightly above zero, and explain why it has the famous ``S'' shape.
\end{exercise}

\begin{exercise}
\label{ch36:ex-superposition1}
The equation $\dv{y}{t} = y^2$ is nonlinear. Show that $y_1 = -1/t$ is a solution (for $t>0$), but that $2y_1 = -2/t$ is \emph{not} a solution. Why does this not violate any principle? Contrast with what would happen for a linear homogeneous equation, and connect your answer to the discussion of superposition in Section~\ref{ch36:sec-classify}.
\end{exercise}
